% Options for packages loaded elsewhere
\PassOptionsToPackage{unicode}{hyperref}
\PassOptionsToPackage{hyphens}{url}
\PassOptionsToPackage{dvipsnames,svgnames,x11names}{xcolor}
\documentclass[
  10pt,
  fontset=fandol]{ctexart}
\usepackage{xcolor}
\usepackage[margin=16mm]{geometry}
\usepackage{amsmath,amssymb}
\setcounter{secnumdepth}{-\maxdimen} % remove section numbering
\usepackage{iftex}
\ifPDFTeX
  \usepackage[T1]{fontenc}
  \usepackage[utf8]{inputenc}
  \usepackage{textcomp} % provide euro and other symbols
\else % if luatex or xetex
  \usepackage{unicode-math} % this also loads fontspec
  \defaultfontfeatures{Scale=MatchLowercase}
  \defaultfontfeatures[\rmfamily]{Ligatures=TeX,Scale=1}
\fi
\usepackage{lmodern}
\ifPDFTeX\else
  % xetex/luatex font selection
\fi
% Use upquote if available, for straight quotes in verbatim environments
\IfFileExists{upquote.sty}{\usepackage{upquote}}{}
\IfFileExists{microtype.sty}{% use microtype if available
  \usepackage[]{microtype}
  \UseMicrotypeSet[protrusion]{basicmath} % disable protrusion for tt fonts
}{}
\makeatletter
\@ifundefined{KOMAClassName}{% if non-KOMA class
  \IfFileExists{parskip.sty}{%
    \usepackage{parskip}
  }{% else
    \setlength{\parindent}{0pt}
    \setlength{\parskip}{6pt plus 2pt minus 1pt}}
}{% if KOMA class
  \KOMAoptions{parskip=half}}
\makeatother
\usepackage{longtable,booktabs,array}
\newcounter{none} % for unnumbered tables
\usepackage{calc} % for calculating minipage widths
% Correct order of tables after \paragraph or \subparagraph
\usepackage{etoolbox}
\makeatletter
\patchcmd\longtable{\par}{\if@noskipsec\mbox{}\fi\par}{}{}
\makeatother
% Allow footnotes in longtable head/foot
\IfFileExists{footnotehyper.sty}{\usepackage{footnotehyper}}{\usepackage{footnote}}
\makesavenoteenv{longtable}
\setlength{\emergencystretch}{3em} % prevent overfull lines
\providecommand{\tightlist}{%
  \setlength{\itemsep}{0pt}\setlength{\parskip}{0pt}}
\usepackage{amsmath,amssymb,mathtools,needspace}
\xeCJKDeclareCharClass{CJK}{"2032,"0400->"04FF,"2160->"216B,"2460->"2473,"25A0,"25C7,"25CB,"25CF}
\IfFontExistsTF{Arial Unicode MS}{\xeCJKsetup{AutoFallBack=true}\setCJKfallbackfamilyfont{\CJKrmdefault}{Arial Unicode MS}}{}
\setcounter{secnumdepth}{0}
\setlength{\emergencystretch}{2em}
\allowdisplaybreaks
\usepackage{bookmark}
\IfFileExists{xurl.sty}{\usepackage{xurl}}{} % add URL line breaks if available
\urlstyle{same}
\hypersetup{
  pdftitle={梯形法的递推化},
  colorlinks=true,
  linkcolor={Maroon},
  filecolor={Maroon},
  citecolor={Blue},
  urlcolor={Blue},
  pdfcreator={LaTeX via pandoc}}

\title{梯形法的递推化}
\author{}
\date{}

\begin{document}
\maketitle

\subsubsection{4.3.1 梯形法的递推化}\label{ux68afux5f62ux6cd5ux7684ux9012ux63a8ux5316}

上一节介绍的复化求积方法对提高精度是行之有效的，但在使用求积公式之前必须给出合适的步长。步长太大，精度难以保证；步长太小，又会导致计算量的增加。而事先给出一个恰当的步长往往是困难的。

实际计算中常常采用变步长的计算方案，即在步长逐次分半（即步长二分）的过程中，反复利用复化求积公式进行计算，直至所求得的积分值满足精度要求为止。

下面，我们在变步长的过程中探讨梯形法的计算规律。

设将求积区间 \([a,b]\) 分成 \(n\) 等份，则一共有 \(n+1\) 个分点，按梯形公式（4.2.9）计算积分值 \(T_n\)，需要提供 \(n+1\) 个函数值。如果将求积区间再二分一次，则分点增至 \(2n+1\) 个。将二分前后两个积分值联系起来加以考察，注意到每个子区间 \([x_k,x_{k+1}]\) 经过二分只增加了一个分点 \(x_{k+\frac12}=\dfrac12(x_k+x_{k+1})\)，用复化梯形公式求得该子区间上的积分值为

\[
\frac h4[f(x_k)+2f(x_{k+\frac12})+f(x_{k+1})].
\]

注意，这里 \(h=\dfrac{b-a}{n}\) 代表二分前的步长。将每个子区间上的积分值相加，得

\[
T_{2n}=\frac h4\sum_{k=0}^{n-1}[f(x_k)+f(x_{k+1})]+\frac h2\sum_{k=0}^{n-1}f(x_{k+\frac12}),
\]

从而利用式（4.2.9）可导出下列递推公式：

\[
T_{2n}=\frac12T_n+\frac h2\sum_{k=0}^{n-1}f(x_{k+\frac12}).
\tag{4.3.1}
\]

\textbf{例 4.2} 计算积分值 \(I=\displaystyle\int_0^1\frac{\sin x}{x}\,\mathrm{d}x\)。

\textbf{解} 先对整个区间 \([0,1]\) 使用梯形公式。对于函数 \(f(x)=\dfrac{\sin x}{x}\)，它在 \(x=0\) 的值定义为 \(f(0)=1\)，而 \(f(1)=0.841\,470\,9\)，据梯形公式计算得

\[
T_1=\frac12[f(0)+f(1)]=0.920\,735\,5.
\]

然后将区间二等分，再求出中点的函数值 \(f(1/2)=0.958\,851\,0\)，从而利用递推公式（4.3.1），有

\href{../../source-images/pdf-102.jpeg}{原页图像}

\[
T_2=\frac12T_1+\frac12f\left(\frac12\right)=0.939\,793\,3.
\]

进一步二分求积区间，并计算新分点上的函数值，得

\[
f(1/4)=0.989\,615\,8,\quad f(3/4)=0.908\,851\,6.
\]

再利用式（4.3.1），有

\[
T_4=\frac12T_2+\frac14\left[f\left(\frac14\right)+f\left(\frac34\right)\right]=0.944\,513\,5.
\]

这样不断二分下去，计算结果如表 4.3 所示（表中 \(k\) 代表二分次数，区间等份数 \(n=2^k\)）。

表 4.3

{\def\LTcaptype{none} % do not increment counter
\begin{longtable}[]{@{}
  >{\raggedright\arraybackslash}p{(\linewidth - 10\tabcolsep) * \real{0.1667}}
  >{\raggedright\arraybackslash}p{(\linewidth - 10\tabcolsep) * \real{0.1667}}
  >{\raggedright\arraybackslash}p{(\linewidth - 10\tabcolsep) * \real{0.1667}}
  >{\raggedright\arraybackslash}p{(\linewidth - 10\tabcolsep) * \real{0.1667}}
  >{\raggedright\arraybackslash}p{(\linewidth - 10\tabcolsep) * \real{0.1667}}
  >{\raggedright\arraybackslash}p{(\linewidth - 10\tabcolsep) * \real{0.1667}}@{}}
\toprule\noalign{}
\begin{minipage}[b]{\linewidth}\raggedright
\(k\)
\end{minipage} & \begin{minipage}[b]{\linewidth}\raggedright
1
\end{minipage} & \begin{minipage}[b]{\linewidth}\raggedright
2
\end{minipage} & \begin{minipage}[b]{\linewidth}\raggedright
3
\end{minipage} & \begin{minipage}[b]{\linewidth}\raggedright
4
\end{minipage} & \begin{minipage}[b]{\linewidth}\raggedright
5
\end{minipage} \\
\midrule\noalign{}
\endhead
\bottomrule\noalign{}
\endlastfoot
\(T_n\) & \(0.939\,793\,3\) & \(0.944\,513\,5\) & \(9.945\,690\,9\) & \(0.945\,985\,0\) & \(0.946\,059\,6\) \\
\end{longtable}
}

{\def\LTcaptype{none} % do not increment counter
\begin{longtable}[]{@{}
  >{\raggedright\arraybackslash}p{(\linewidth - 10\tabcolsep) * \real{0.1667}}
  >{\raggedright\arraybackslash}p{(\linewidth - 10\tabcolsep) * \real{0.1667}}
  >{\raggedright\arraybackslash}p{(\linewidth - 10\tabcolsep) * \real{0.1667}}
  >{\raggedright\arraybackslash}p{(\linewidth - 10\tabcolsep) * \real{0.1667}}
  >{\raggedright\arraybackslash}p{(\linewidth - 10\tabcolsep) * \real{0.1667}}
  >{\raggedright\arraybackslash}p{(\linewidth - 10\tabcolsep) * \real{0.1667}}@{}}
\toprule\noalign{}
\begin{minipage}[b]{\linewidth}\raggedright
\(k\)
\end{minipage} & \begin{minipage}[b]{\linewidth}\raggedright
6
\end{minipage} & \begin{minipage}[b]{\linewidth}\raggedright
7
\end{minipage} & \begin{minipage}[b]{\linewidth}\raggedright
8
\end{minipage} & \begin{minipage}[b]{\linewidth}\raggedright
9
\end{minipage} & \begin{minipage}[b]{\linewidth}\raggedright
10
\end{minipage} \\
\midrule\noalign{}
\endhead
\bottomrule\noalign{}
\endlastfoot
\(T_n\) & \(0.946\,076\,9\) & \(0.946\,081\,5\) & \(0.946\,082\,7\) & \(0.946\,083\,0\) & \(0.946\,083\,1\) \\
\end{longtable}
}

积分 \(I\) 的准确值为 \(0.946\,083\,1\)，用变步长方法二分 9 次得到了这个结果。

\begin{quote}
校注（agent 补充，CH04A-E03）：表 4.3 的 \(k=3\) 单元格原扫描确实印为 \(9.945\,690\,9\)，这里按原表保留。由例 4.1 的 \(T_8\) 和本页下段的 \(T_8\) 可知首位应为 \(0\)，即 \(0.945\,690\,9\)，这是原书排印错误。

校注（agent 补充，CH04A-E04）：本页说明 \(k\) 为二分次数，表中 \(k=9\) 为 \(0.946\,083\,0\)，\(k=10\) 才列出 \(0.946\,083\,1\)。故表后``二分 9 次得到了这个结果''与本页表格、次数定义不一致；原句已保留。

校注（agent 补充，CH04A-E07）：表 4.3 的 \(k=5\) 单元格原扫描印为 \(0.946\,059\,6\)，这里按原表保留。对本例 \(f(x)=\sin x/x\)（\(f(0)=1\)）在 \([0,1]\) 上取 \(n=32\) 独立高精度计算复化梯形值，得 \(T_{32}=0.946\,058\,560\,962\,768\,067\ldots\)，保留七位小数应为 \(0.946\,058\,6\)；原表数值有误。
\end{quote}

\subsection{本节覆盖原页的校注记录}\label{ux672cux8282ux8986ux76d6ux539fux9875ux7684ux6821ux6ce8ux8bb0ux5f55}

下列记录按原页关联，含同页相邻小节的校注；计算时同时读取上文校注和完整原页。

\begin{itemize}
\tightlist
\item
  \texttt{CH04A-E03}：原书 PDF 页 102
\item
  \texttt{CH04A-E04}：原书 PDF 页 102
\item
  \texttt{CH04A-E07}：原书 PDF 页 102
\end{itemize}

\end{document}
